\label{der:3-3}
记旧、新基矩阵为$B,C$且$C=BP$。同一向量分别写为$\vx=B\xi=C\eta=BP\eta$。因为旧基列无关，$B(\xi-P\eta)=0$推出$\xi=P\eta$。若$Pz=0$，则$Cz=BPz=0$，新基无关使$z=0$；$P$为方阵，故可逆，得到$\eta=P\inv\xi$。

为说明任意两组基个数相同，设长度为$m,n$。把第一组基用第二组表示，其系数列必须无关，否则对应的非零系数组合会使第一组相关；这些系数列处在$\F^n$，故$m\le n$。互换两组得$n\le m$，所以相等。若$n$维空间中有$n$个无关向量却不生成全空间，可加入一个不在其生成空间中的向量，得到$n+1$个无关向量，矛盾。若$n$个向量生成全空间却相关，删去一个可由其余表示的向量，剩下至多$n-1$个仍生成，会与空间维数$n$矛盾。因此两种单独条件都足以构成基。
